
\section{Preliminaries}\label{sec:prelims}

We use \textbf{boldface font} to denote random variables.
For two complex numbers $\alpha, \beta \in \C$, we write
$\alpha \approx_{\eps} \beta$ if
\begin{equation*}
|\alpha - \beta| \leq \eps.
\end{equation*}
We note the following triangle inequality for numbers, which we will use repeatedly:
\begin{equation}\label{eq:triangle-inequality-for-numbers}
\text{if $\alpha \approx_{\eps} \beta$ and $\beta \approx_{\delta} \gamma$, then $\alpha \approx_{\eps+\delta}\gamma$.}
\end{equation}

\subsection{Finite fields}

A \emph{finite field} is a field with a finite number of elements.
There is a unique finite field~$\F_q$ of~$q$ elements for each prime power $q = p^t$,
and there are no other finite fields.
We write $\omega$ for the $p$-th root of unity $\omega = e^{2 \pi i/p}$.

\begin{definition}[Finite field trace]
The \emph{finite field trace} is the function $\tr:\F_q \rightarrow \F_p$ defined as
\begin{equation*}
\tr[x] = \sum_{\ell=0}^{t-1} x^{p^{\ell}}.
\end{equation*}
\end{definition}

\begin{proposition}
\label{prop:fourier-fact-scalar}
Let $a \in \F_q$. Then
\begin{equation*}
\E_{\bx \sim \F_q} \omega^{\tr[\bx \cdot a]}
= \left\{\begin{array}{rl}
	1 & \text{if $a = 0$},\\
	0 & \text{otherwise}.
	\end{array}\right.
\end{equation*}
\end{proposition}
\begin{proof}
If $a = 0$, then $\tr[x \cdot a] = 0$ for all $x \in \F_q$.
As a result,
\begin{equation*}
\E_{\bx \sim \F_q} \omega^{\tr[\bx \cdot a]}
= \E_{\bx \sim \F_q} \omega^{0}
= \E_{\bx \sim \F_q} 1
= 1.
\end{equation*}
On the other hand, if $a \neq 0$, then there exists a $y \in \F_q$
such that $\tr[a \cdot y] \neq 0$. As a result,
\begin{equation*}
C
:= \E_{\bx \sim \F_q} \omega^{\tr[\bx \cdot a]}
= \E_{\bx \sim \F_q} \omega^{\tr[(\bx+y) \cdot a]}
= \E_{\bx \sim \F_q} (\omega^{\tr[\bx \cdot a]} \cdot \omega^{\tr[y \cdot a]})
= C \cdot \omega^{\tr[y\cdot a]}.
\end{equation*}
But because $\tr[y \cdot a] \neq 0$,
$\omega^{\tr[y\cdot a]} \neq 1$.
This implies that $C = 0$.
\end{proof}


\begin{proposition}
\label{prop:fourier-fact-vector}
Let $v \in \F_q^m$. Then
\begin{equation*}
\E_{\bu \sim \F_q^m} \omega^{\tr[\bu \cdot v]}
= \left\{\begin{array}{rl}
	1 & \text{if $v = 0$},\\
	0 & \text{otherwise}.
	\end{array}\right.
\end{equation*}
\end{proposition}
\begin{proof}
By the linearity of the trace,
\begin{equation*}
\E_{\bu \sim \F_q^m} \omega^{\tr[\bu \cdot v]}
= \E_{\bu \sim \F_q^m} (\omega^{\tr[\bu_1 \cdot v_1]} \cdots \omega^{\tr[\bu_m \cdot v_m]})
= \prod_{i=1}^m (\omega^{\tr[\bu_i \cdot v_i]}),
\end{equation*}
which by \Cref{prop:fourier-fact-scalar} is $1$ if $v_i = 0$ for all $i$ and $0$ otherwise.
\end{proof}



\subsection{Polynomials over finite fields}

\begin{definition}[Polynomials of low individual degree]
Let $q$ be a prime power,
and let $m$ and $d$ be nonnegative integers.
We define $\polyfunc{m}{q}{d}$ to be the set of polynomials in $\F_q^m$
with individual degree~$d$.
\end{definition}

\begin{remark}
Note that as defined in \Cref{sec:intro},
in this work a polynomial with individual degree~$d$ is one in which the degree~$d_i$ of each coordinate~$i$
is \emph{at most~$d$}.
This allows us to say ``individual degree~$d$'' rather than the wordier ``individual degree at most~$d$''.
For example, under this definition, $\polyfunc{m}{q}{d}$ is contained in $\polyfunc{m}{q}{d+1}$ for each~$d$.
\end{remark}

The most important fact about low-degree polynomials is that they have large distance from each other.
This is shown by the following lemma.

\begin{lemma}[Schwartz-Zippel lemma~\cite{Sch80,Zip79}]\label{lem:schwartz-zippel-total-degree}
Let $g, h:\F_q^m \rightarrow \F_q$ be two distinct polynomials of total degree~$d$.
Then
\begin{equation*}
\Pr_{\bx \sim \F_q^m}[g(\bx) = h(\bx)] \leq \frac{d}{q}.
\end{equation*}
\end{lemma}

Since any polynomial with individual degree~$d$ has total degree~$md$,
\Cref{lem:schwartz-zippel-total-degree} implies the following corollary.

\begin{corollary}[Schwartz-Zippel for individual degree]
Let $g, h \in \polyfunc{m}{q}{d}$ be distinct. Then
\begin{equation*}
\Pr_{\bx \sim \F_q^m}[g(\bx) = h(\bx)] \leq \frac{md}{q}.
\end{equation*}
\end{corollary}

\subsection{Measurements}

\begin{definition}[Measurements and sub-measurements]
Let $\calH$ be a Hilbert space and $\calA$ be a set of outcomes.
A \emph{sub-measurement}
is a set of Hermitian, positive-semidefinite matrices
$A = \{A_a\}_{a \in \calA}$
acting on~$\calH$
such that $\sum_a A_a \leq I$.
The sub-measurement is \emph{projective} if $(A_a)^2 = A_a$ for each~$a$.
It is a \emph{measurement}
if it satisfies the stronger condition $\sum_a A_a = I$.
\end{definition}

An important class of sub-measurements
are those that output polynomials~$g$ of individual degree~$d$.
These are defined as follows.

\begin{definition}[Low-degree polynomial measurements]
We write $\polysub{m}{q}{d}$ for the set of sub-measurements $G=\{G_g\}$
with outcomes $g \in \polyfunc{m}{q}{d}$.
We write $\polymeas{m}{q}{d}$ for the subset of $\polysub{m}{q}{d}$
containing only those $G = \{G_g\}$ which are measurements.
\end{definition}

\begin{definition}[Post-processing measurements]\label{def:post-processing}
Let $A = \{A_a\}_{a \in \calA}$ be a set of matrices,
and let $f:\calA \rightarrow \calB$ be a function.
Then for each $b \in \calB$,
we define the matrix
\begin{equation*}
A_{[f(a)=b]} = \sum_{a:f(a)=b} A_a.
\end{equation*}
\end{definition}

\begin{remark}
We note that \Cref{def:post-processing}
agrees with the notation for post-processing measurements
used in \cite{NW19},
but it disagrees with the notation used in \cite{JNV+20}.
That work uses the notation ``$A_{[f(\cdot)=b]}$'' rather
than the notation ``$A_{[f(a)=b]}$'' that we use in this work. 
\end{remark}

An easy-to-prove fact is that measurements remain measurements after post-processing their outcomes.

\begin{proposition}
Let $A = \{A_a\}_{a \in \calA}$ be a set of matrices,
and let $f:\calA \rightarrow \calB$ be a function.
Then
\begin{equation*}
\sum_a A_a = \sum_b A_{[f(a)=b]}.
\end{equation*}
Thus, if~$\{A_a\}$ is a sub-measurement (respectively, measurement),
then~$\{A_{[f(a)=b]}\}$ is also a sub-measurement (respectively, measurement).
\end{proposition}

\begin{remark}
We note that there is some ambiguity in the notation $A_{[f(a)=b]}$
because it requires one to know which of~$f$ or~$a$ is the measurement outcome
and which is the function applied to it.
For example, given a sub-measurement $G = \{G_g\} \in \polysub{m}{q}{d}$,
we will often consider evaluating its outputs at a point $u \in \F_q^m$.
This entails looking at the sub-measurement
\begin{equation*}
\{G_{[g(u)=a]}\}_{a \in \F_q},
\quad
\text{where}
\quad
G_{[g(u)=a]} = \sum_{g:g(u) = a} G_g.
\end{equation*}
Here, $g$ is the measurement outcome of~$G$,
and the function being applied maps it to its evaluation on the point~$u$, i.e.\ $g(u)$.
In general, it should always be clear from context
what the measurement outcome and the function being applied to it are.
\end{remark}

\begin{notation}[Measurements indexed by questions]
We will frequently encounter sets of sub-measurements $A^x= \{A^x_{a}\}$ indexed by elements~$x$ from a set of ``questions" $\calX$. We will write $A = \{A^x_a\}$ for this set. 
We will typically refer to this set~$A$ as a sub-measurement,
and we will refer to it as a measurement if each~$A^x$ is a measurement.
In addition, we will refer to it as projective if each~$A^x$ is projective.
\end{notation}

\begin{notation}[Complete part of sub-measurement]
Given a sub-measurement~$A = \{A_a\}$, we will write $A = \sum_a A_a$.
Similarly, given a sub-measurement~$A = \{A^x_a\}$, we will write $A^x = \sum_a A^x_a$ and $A = \E_{\bx} A^{\bx}$.
\end{notation}

The complete part of a sub-measurement
contrasts with its \emph{incomplete part},
which is the matrix $I - A^x$.
We will sometimes
view~$A$ as a measurement
by throwing in its incomplete part as an additional POVM element.
This is formalized in the following definition.


\begin{definition}[Completing a sub-measurement]\label{def:measurement-completion}
Let $A = \{A^x_a\}_{a \in \calA}$ be a sub-measurement.
We define the \emph{completion of~$A$},
denoted $\mathrm{completion}(A)$,
to be the measurement $\widehat{A} = \{\widehat{A}^x_a\}$
with outcome set $\widehat{\calA} = \calA \cup \{\bot\}$ 
such that for each $a \in \widehat{\calA}$,
\begin{equation*}
A^x_a
= \left\{\begin{array}{cl}
		A^x_a & \text{if $a \in \calA$,}\\
		I - A^x & \text{if $a = \bot$}.
		\end{array}\right.
\end{equation*}
\end{definition}

\subsection{Comparing measurements}
\label{sec:comparing-measurements}

An central problem in this paper is recognizing when two measurements are close to each other. 
We will survey two methods of doing so,
using the \emph{consistency} and the \emph{state dependent distance}.
This section largely mirrors Sections 4.4 and 4.5 of \cite{NW19},
which prove numerous properties of these two distances.
However, we are unable to cite their results directly
because they are mostly stated and proven only for measurements,
whereas we will need to apply them to sub-measurements as well.

\subsubsection{Consistency between measurements}

The most basic notion of similarity between two measurements is given by their \emph{consistency}.

\begin{definition}[Consistency]\label{def:simeq}
Let $\ket{\psi}$ be a state in $\calH_{\mathrm{A}} \ot \calH_{\mathrm{B}}$.
Let $A = \{A^x_a\}$ be a sub-measurement acting on $\calH_{\mathrm{A}}$
and $B = \{B^x_a\}$ be a sub-measurement acting on $\calH_{\mathrm{B}}$.
Finally, let $\calD$ be a distribution on the question set~$\calX$.
Then we say that
 \begin{equation*}
 A^x_a \ot I \simeq_{\delta} I \ot B^x_a
 \end{equation*}
  on state $\ket{\psi}$ and distribution $\calD$ if
  \begin{equation}\label{eq:no-big-Oh}
  \E_{\bx \sim \calD} \sum_{a \neq b} \bra{\psi} A^{\bx}_a \ot B^{\bx}_b
    \ket{\psi} \leq \delta.
    \end{equation}
\end{definition}

This is simply the probability that the provers receive different outcomes
when they measure with~$A$ and~$B$,
assuming the sub-measurements do return an outcome.

\begin{notation}[Simplifying notation]
Because the state $\ket{\psi}$ and distribution $\calD$ are typically clear from context,
we will often write ``$ A^x_a \ot I \simeq_{\delta} I \ot B^x_a$"
as shorthand for ``$ A^x_a \ot I \simeq_{\delta} I \ot B^x_a$ on state $\ket{\psi}$ and distribution $\calD$".
In the case when $\calD$ is not clear by context,
we might specify it implicitly in terms of a random variable~$\bx$ distributed according to~$\calD$.
For example, suppose the distribution~$\calD$ is supposed to be uniform on the question set $\F_q$.
Then we might say ``on average over $\bx \sim \F_q$,
\begin{equation*}
A^x_a \ot I \simeq_{\delta} I \ot B^x_a"
\end{equation*}
as shorthand for ``$ A^x_a \ot I \simeq_{\delta} I \ot B^x_a$ on distribution $\calD$".
\end{notation}

We note that there are two differences between the definition of consistency in \Cref{def:simeq}
and the original definition of consistency given in \cite[Definition 4.11]{NW19}.
The first of these is that the~\cite{NW19} definition
allows the right-hand side of \Cref{eq:no-big-Oh} to be $O(\delta)$
rather than strictly~$\delta$.
The benefit of this is that they do not need to keep track of constant prefactors in their proofs;
we have elected to use this more concrete definition
to make our proofs more easily verifiable, at the expense of tracking these constant prefactors.
The second difference is that is that they define their consistency as the quantity
\begin{equation*}
1 - \E_{\bx} \sum_a \bra{\psi} A^{\bx}_a \ot B^{\bx}_a \ket{\psi}.
\end{equation*}
As we show below in \Cref{prop:simeq-for-measurements},
this agrees with \Cref{def:simeq}
when~$A$ and~$B$ are both measurements.
However, these two definitions disagree
when~$A$ and~$B$ are \emph{sub}-measurements,
which is a case we will frequently encounter throughout this paper.

\begin{proposition}[Consistency for measurements]\label{prop:simeq-for-measurements}
Let $A = \{A^x_a\}$ and $B = \{B^x_a\}$ be two measurements.
Then
\begin{equation*}
A^x_a \ot I \simeq_{\delta} I \ot B^x_a
\end{equation*}
if and only if
\begin{equation*}
\E_{\bx} \sum_a \bra{\psi} A^{\bx}_a \ot B^{\bx}_a \ket{\psi} \geq 1 - \delta.
\end{equation*}
\end{proposition}
\begin{proof}
We compute
\begin{align*}
&\E_{\bx} \sum_{a \neq b}\bra{\psi} A^{\bx}_a \ot B^{\bx}_b \ket{\psi}\\
=~& \E_{\bx} \sum_b \bra{\psi} \Big(\sum_{a \neq b} A^{\bx}_a\Big) \ot B^{\bx}_b\ket{\psi}\\
=~&\E_{\bx} \sum_b \bra{\psi} (I - A^{\bx}_b) \ot B^{\bx}_b\ket{\psi}
	\tag{because~$A$ is a measurement}\\
=~&\E_{\bx} \sum_b \bra{\psi} I \ot B^{\bx}_b\ket{\psi}
	- \E_{\bx} \sum_b \bra{\psi} A^{\bx}_b \ot B^{\bx}_b\ket{\psi}\\
=~&1
	- \E_{\bx} \sum_b \bra{\psi} A^{\bx}_b \ot B^{\bx}_b\ket{\psi}. \tag{because~$B$ is a measurement}
\end{align*}
Hence, the first expression is at most~$\delta$ if and only if the last one is as well,
and we are done.
\end{proof}

\subsubsection{The state-dependent distance}

Suppose we have three measurements $\{A^x_a\}$, $\{B^x_a\}$, and $\{C^x_a\}$,
and we know that
\begin{equation*}
A^x_a \ot I \simeq_{\eps} I \ot C^x_a.
\end{equation*}
What property of~$A$ and~$B$ allows us to conclude that
\begin{equation}\label{eq:can-we-use-approx-delta-to-derive-this}
B^x_a \ot I \simeq_{\eps} I \ot C^x_a?
\end{equation}
The answer is provided by the \emph{state-dependent distance}.

\begin{definition}[The state-dependent distance]\label{def:approx_delta}
Let $\ket{\psi}$ be a state in $\calH$.
Let $A = \{A^x_a\}$ and $B = \{B^x_a\}$ be
sets of matrices acting on $\calH$.
Finally, let $\calD$ be a distribution on the question set~$\calX$.
Then we say that
  \begin{equation*}
  A^x_a \approx_{\delta}  B^x_a
  \end{equation*}
 on state $\ket{\psi}$ and distribution $\calD$ if
  \begin{equation*}
  \E_{\bx \sim \calD} \sum_{a} \Vert (A^{\bx}_a - B^{\bx}_a) \ket{\psi} \Vert^2 \leq \delta.
    \end{equation*}
\end{definition}

We note one odd feature of \Cref{def:approx_delta} in comparison to the consistency,
which is that $\ket{\psi}$ is not assumed to have a bipartition $\calH_{\mathrm{A}} \ot \calH_{\mathrm{B}}$
in which~$A$ and~$B$ are applied on opposite sides.
We will address this in \Cref{sec:consistency-from-state-dependent-distance} below.
Before doing so, we will answer our question above,
even in the case when~$C$ is allowed to be a sub-measurement. 

\begin{proposition}[Transfering ``$\simeq$" using ``$\approx$"]
Let $\{A^x_a\}$ and $\{B^x_a\}$ be measurements,
and let $\{C^x_a\}$ be a sub-measurement.
Suppose that $A^x_a \ot I \simeq_{\delta} I \ot C^x_a$
and $A^x_a \ot I \approx_{\eps} B^x_a \ot I$.
Then $B^x_a \ot I \simeq_{\delta + \sqrt{\eps}} I \ot C^x_a$.
  \label{prop:triangle-sub}
\end{proposition}


\begin{proof}
  First, we can
  rewrite the consistency between $A$ and $C$ as
  \begin{align*}
    \delta \geq \E_{\bx} \sum_{a \neq b} \bra{\psi} A^{\bx}_a \ot C^{\bx}_b
    \ket{\psi} 
    & = \E_{\bx} \sum_{a} \bra{\psi} A^{\bx}_a \ot (C^{\bx} - C^{\bx}_a) \ket{\psi} \\
    & =  \E_{\bx} \sum_{a} \bra{\psi} A^{\bx}_a \ot C^{\bx} \ket{\psi}
    			- \E_{\bx} \sum_{a} \bra{\psi} A^{\bx}_a \ot C^{\bx}_a \ket{\psi} \\
    & =  \E_{\bx} \bra{\psi} I \ot C^{\bx} \ket{\psi}
    			- \E_{\bx} \sum_{a} \bra{\psi} A^{\bx}_a \ot C^{\bx}_a \ket{\psi} \tag{because~$A$ is a measurement}\\
    & = \bra{\psi} I \ot C \ket{\psi} - \E_{\bx} \sum_a
                 \bra{\psi} A^{\bx}_a \ot C^{\bx}_a \ket{\psi}.
  \end{align*}
  Likewise, we can rewrite the
  inconsistency between $B$ and $C$ as
  \begin{align*}
    \E_{\bx} \sum_{a \neq b} \bra{\psi} B^{\bx}_a \ot C^{\bx}_b
    \ket{\psi} &= \bra{\psi} I \ot C \ket{\psi} - \E_{\bx} \sum_a
                 \bra{\psi} B^{\bx}_a \ot C^{\bx}_a \ket{\psi}.
  \end{align*}
  We want to show that the inconsistency between~$B$ and~$C$ is close to the inconsistency between $A$
  and $B$. In particular, we claim that
  \begin{equation*}
  \E_{\bx} \sum_{a \neq b} \bra{\psi} A^{\bx}_a \ot C^{\bx}_b
    \ket{\psi}
    \approx_{\sqrt{\eps}} \E_{\bx} \sum_{a \neq b} \bra{\psi} B^{\bx}_a \ot C^{\bx}_b
    \ket{\psi}.
  \end{equation*}
  To show this, we bound the magnitude of the difference using Cauchy-Schwarz.
  \begin{align*}
   & \Big| \E_{\bx} \sum_{a } \bra{\psi} (A^{\bx}_a - B^{\bx}_a) \ot
                     C^{\bx}_a\ket{\psi} \Big| \nonumber
    \\
    \leq~& \sqrt{ \E_{\bx} \sum_{a} \bra{\psi}
                            (A^{\bx}_a - B^{\bx}_a)^2 \ot I \ket{\psi}}
                            \cdot \sqrt{\E_{\bx} \sum_{a} \bra{\psi} I \ot (C^{\bx}_a)^2 \ket{\psi}} \\
    \leq~& \sqrt{\eps} \cdot 1. \tag{because~$C$ is a sub-measurement}
  \end{align*}
  This completes the proof.
\end{proof}


Next, we show that in the case of measurements,
the state-dependent distance is a weakening of the consistency.

\begin{proposition}[``$\simeq$" implies ``$\approx$" for measurements]\label{prop:simeq-to-approx}
Let $A = \{A^x_a\}$ and $B = \{B^x_a\}$ be two measurements such that
\begin{equation*}
A^x_a \ot I \simeq_{\delta} I \ot B^x_a.
\end{equation*}
Then
\begin{equation*}
A^x_a \ot I \approx_{2\delta} I \ot B^x_a.
\end{equation*}
This is an ``if and only if'' if~$A$ and~$B$ are both projective measurements.
\end{proposition}
\begin{proof}
Our goal is to bound
\begin{align*}
&\E_{\bx} \sum_a\Vert (A^{\bx}_a \ot I - I \ot B^{\bx}_a) \ket{\psi} \Vert^2\\
=~&\E_{\bx} \sum_a \bra{\psi} (A^{\bx}_a \ot I - I \ot B^{\bx}_a)^2 \ket{\psi}\\
=~&\E_{\bx} \sum_a \bra{\psi} (A^{\bx}_a)^2 \ot I \ket{\psi}
	+\E_{\bx} \sum_a \bra{\psi} I \ot (B^{\bx}_a)^2 \ket{\psi}
	- 2 \E_{\bx} \sum_a \bra{\psi} A^{\bx}_a \ot B^{\bx}_a \ket{\psi}\\
\leq~&\E_{\bx} \sum_a \bra{\psi} A^{\bx}_a \ot I \ket{\psi}
	+\E_{\bx} \sum_a \bra{\psi} I \ot B^{\bx}_a \ket{\psi}
	- 2 \E_{\bx} \sum_a \bra{\psi} A^{\bx}_a \ot B^{\bx}_a \ket{\psi}\\
=~&2
	- 2 \E_{\bx} \sum_a \bra{\psi} A^{\bx}_a \ot B^{\bx}_a \ket{\psi} \tag{because~$A$ and~$B$ are measurements}\\
\leq~& 2 - 2 \cdot(1-\delta) \tag{by~\Cref{prop:simeq-for-measurements} and the fact that~$A$ and~$B$ are measurements}\\
=~&2 \delta.
\end{align*}
This completes the proof.
When  $A$ and~$B$ are projective, ``if and only if'' 
follows from the first inequality becoming an equality.
\end{proof}


\begin{remark}
We note that \Cref{prop:simeq-to-approx} certainly does \emph{not} hold for sub-measurements.
For example, if $A^x_a = 0$ for all~$a$,
then $A_a^x \ot I \simeq_{0} I \ot B_a^x$,
but
\begin{equation*}
\E_{\bx} \sum_a \Vert (A^{\bx}_a \ot I - I \ot B^{\bx}_a)\ket{\psi} \Vert^2
= \E_{\bx} \sum_a \Vert (I \ot B^{\bx}_a)\ket{\psi} \Vert^2,
\end{equation*}
which is nonzero unless $(I \ot B^x_a) \ket{\psi} = 0$ for all~$x$ and~$a$.
\end{remark}


\subsubsection{Deriving consistency relations from the state-dependent distance}
\label{sec:consistency-from-state-dependent-distance}

\Cref{prop:triangle-sub} shows the purpose of the state-dependent distance,
which is to derive new consistency relations from old ones.
In addition, its proof uses a strategy
which will recur frequently throughout this paper.
In this strategy, we would like to demonstrate a sequence of 
expressions in which each expression is close to the previous one:
\begin{equation*}
\E_{\bx} \sum_a \bra{\psi} (A_0)^{\bx}_a (B_0)^{\bx}_a \ket{\psi}
\approx_{\eps_1} \E_{\bx} \sum_a \bra{\psi} (A_1)^{\bx}_a (B_1)^{\bx}_a \ket{\psi}
\approx_{\eps_2} \cdots
\approx_{\eps_t} \E_{\bx} \sum_a \bra{\psi} (A_t)^{\bx}_a (B_t)^{\bx}_a \ket{\psi}.
\end{equation*}
By the triangle inequality, we can therefore conclude that the $0$-th expression
is close to the $t$-th expression.
To show that the $i$-th quantity is close to the $(i+1)$-st,
we will typically arrange for $A_{i} = A_{i+1}$,
and we will swap out $B_i$ for $B_{i+1}$
using an approximation relation such as $(B_i)^x_a \approx (B_{i+1})^x_a$,
with the help of the Cauchy-Schwarz inequality
(or the same might occur with the roles of~$A$ and~$B$ reversed). 
Even if $A_0$ and $B_0$ can be written as local measurements applied to either side of a bipartition,
e.g.\ $(A_0)^x_a = A^x_a \ot I$ and $(B_0)^x_a = I \ot B^x_a$,
and likewise for $A_t$ and $B_t$,
the intermediate steps may feature matrices which do not decompose nicely across a bipartition.
This is why \Cref{def:approx_delta} is phrased so broadly,
with no mention of a bipartition.

In general, the $A_i$'s and $B_i$'s encountered in this sequence of steps
may be quite unstructured: for example, not sub-measurements,
and possibly not even Hermitian.
Thus, we are interested in determining which conditions
to place on these measurements are sufficient to carry out this proof strategy.
The following proposition gives a broad condition under which this can be accomplished.

\begin{proposition}
\label{prop:closeness-of-ip}
	Let $\{A^x_a\}$, $\{B^x_a\}$, and $\{C^x_{a,b} \}$ be matrices.
	Suppose that  $A^x_a \approx_\gamma B^x_a$ and that for all~$x$, $\sum_a (\sum_b C^x_{a,b}) (\sum_b C^x_{a,b})^\dagger  \leq I$. Then
	\begin{equation}
	\E_{\bx} \sum_{a,b} \bra{\psi} C^{\bx}_{a,b} A^{\bx}_a  \ket{\psi} \approx_{\sqrt{\gamma}} \E_{\bx} \sum_{a,b} \bra{\psi} C^{\bx}_{a,b} B^{\bx}_a  \ket{\psi} \,. \label{eq:closeness3}
	\end{equation}
	Similarly, suppose that  $(A^x_a)^\dagger \approx_\gamma (B^x_a)^\dagger$ and that for all~$x$, $\sum_a (\sum_b C^x_{a,b})^\dagger (\sum_b C^x_{a,b}) \leq I$. Then 
	\begin{equation}
	\E_{\bx} \sum_{a,b} \bra{\psi} A^{\bx}_a  C^{\bx}_{a,b} \ket{\psi} \approx_{\sqrt{\gamma}} \E_{\bx} \sum_{a,b} \bra{\psi} B^{\bx}_a  C^{\bx}_{a,b}  \ket{\psi} \,. \label{eq:closeness4}
	\end{equation}	
\end{proposition}
\begin{proof}
	We begin by showing~\Cref{eq:closeness3}.
	\begin{align*}
		&\Big | \E_{\bx} \sum_{a,b} \bra{\psi} C^{\bx}_{a,b} (A^{\bx}_a - B^{\bx}_a) \ket{\psi} \Big | \\
		&= \Big | \E_{\bx} \sum_{a} \bra{\psi} \Big ( \sum_b C^{\bx}_{a,b} \Big ) \cdot (A^{\bx}_a - B^{\bx}_a) \ket{\psi} \Big |\\
		&\leq \Big ( \E_{\bx} \sum_a \bra{\psi} \Big(\sum_b C^{\bx}_{a,b}\Big) \Big(\sum_b C^{\bx}_{a,b}\Big)^\dagger \ket{\psi} \Big )^{1/2} \cdot \Big ( \E_{\bx} \sum_a \bra{\psi} (A^{\bx}_a - B^{\bx}_a)^\dagger  (A^{\bx}_a - B^{\bx}_a) \ket{\psi} \Big)^{1/2} \\
		&\leq \sqrt{\gamma}.
	\end{align*}
	The third line uses Cauchy-Schwarz, and the fourth line uses the assumption $\sum_a (\sum_b C^x_{a,b}) (\sum_b C^x_{a,b})^\dagger  \leq I$ to bound the first factor by $1$ and  the assumption $A^x_a \approx_\gamma B^x_a$ to bound the second factor by $\sqrt{\gamma}$.
	As for \Cref{eq:closeness4}, we want to bound
	\begin{equation*}
	\Big | \E_{\bx} \sum_{a,b} \bra{\psi} (A^{\bx}_a - B^{\bx}_a) C^{\bx}_{a,b}  \ket{\psi} \Big |
	= \Big | \E_{\bx} \sum_{a,b} \bra{\psi} (C^{\bx}_{a,b})^\dagger ((A^{\bx}_a)^\dagger - (B^{\bx}_a)^\dagger) \ket{\psi} \Big |.
      \end{equation*}
      It then follows from \Cref{eq:closeness3} that this is at most $\sqrt{\gamma}$.
    \end{proof}
    
\Cref{prop:closeness-of-ip} is broad enough to capture
almost all of our applications of the state-dependent distance.
Unfortunately, defining the $C^x_{a, b}$ matrices and showing that 
they satisfy the inequality $\sum_a (\sum_b C^x_{a,b}) (\sum_b C^x_{a,b})^\dagger  \leq I$
can be somewhat cumbersome.
As a result, we will usually carry out these Cauchy-Schwarz calculations by hand.
However, we will occasionally use the following proposition
which simplifies \Cref{prop:closeness-of-ip}.

\begin{proposition}\label{prop:easy-approx-from-approx-delta}
Let $A = \{A^x_a\}$, $B = \{B^x_a\}$, and $C = \{C^x_a\}$ be sub-measurements
such that $A^x_a \approx_{\delta} B^x_a$. Then
\begin{equation*}
\E_{\bx} \sum_a \bra{\psi} A^{\bx}_a  C^{\bx}_a \ket{\psi}
\approx_{\sqrt{\delta}} \E_{\bx} \sum_a \bra{\psi} B^{\bx}_a  C^{\bx}_a \ket{\psi}.
\end{equation*}
\end{proposition}
\begin{proof}
To show this, we bound the magnitude of the difference.
\begin{align*}
&\Big| \E_{\bx} \sum_a \bra{\psi} (A^{\bx}_a - B^{\bx}_a) \cdot (C^{\bx}_a) \ket{\psi}\Big|\\
\leq~& \sqrt{\E_{\bx} \sum_a \bra{\psi} (A^{\bx}_a - B^{\bx}_a)^2 \ket{\psi}}
	\cdot \sqrt{\E_{\bx} \sum_a \bra{\psi} (C^{\bx}_a)^2 \ket{\psi}}\\
\leq~& \sqrt{\delta} \cdot \sqrt{1}.
\end{align*}
This completes the proof.
\end{proof}

A proposition similar to \Cref{prop:closeness-of-ip},
but for  ``$\approx$'',
holds as well.
    This is \cite[Fact~$4.20$]{NW19}.
    \begin{proposition}
      \label{prop:cab-approx-delta}
      Let $\{A^x_a\}, \{B^x_a\},$ and $\{C^x_{a,b}\}$ be
      matrices. Suppose that $A^{x}_a \approx_\delta B^{x}_a$ and that
      for all $x$ and $a$, $\sum_b (C^{x}_{a,b})^\dagger
      (C^{x}_{a,b}) \leq I$. Then
      \[ C^{x}_{a,b} A^x_{a} \approx_{\delta} C^{x}_{a,b} B^{x}_a. \]
    \end{proposition}
    \begin{proof}
      The error we wish to bound is
      \begin{align*}
        \E_{\bx} \sum_{a,b} \bra{\psi} 
        (A^{\bx}_a - B^{\bx}_a)^\dagger
        (C^{\bx}_{a,b})^\dagger C^{\bx}_{a,b} (A^{\bx}_a - 
        B^{\bx}_a) \ket{\psi} &\leq \E_{\bx} \sum_{a} \bra{\psi}
                                (A^{\bx}_a - B^{\bx}_a)^\dagger
                                (A^{\bx}_a - B^{\bx}_a) \ket{\psi} \\
                              &\leq \delta. \qedhere
      \end{align*}
    \end{proof}


\subsubsection{Miscellaneous distance properties}

We now state a few miscellaneous properties of our two distances.

\begin{proposition}[Triangle inequality for vectors squared]\label{prop:triangle-inequality-for-vectors-squared}
Let $\ket{\psi_1}, \ldots, \ket{\psi_k}$ be vectors.
Then
\begin{equation*}
\Vert \ket{\psi_1} + \cdots + \ket{\psi_k} \Vert^2 \leq k \cdot(\Vert \ket{\psi_1}\Vert^2 + \cdots + \Vert \ket{\psi_k}\Vert^2).
\end{equation*}
\end{proposition}
\begin{proof}
First, if $x_1, \ldots, x_k \in \R$, then
\begin{equation}\label{eq:prop-for-real-numbers}
(x_1 + \cdots + x_k)^2 
= \sum_{i, j= 1}^k x_i x_j
\leq \sum_{i, j= 1}^k \frac{x_i^2 + x_j^2}{2}
= \sum_{i, j= 1}^k x_i^2
= \sum_{i=1}^k k \cdot x_i^2.
\end{equation}
Next, by the triangle inequality
\begin{align*}
\Vert \ket{\psi_1} + \cdots + \ket{\psi_k} \Vert^2
& = (\Vert \ket{\psi_1} + \cdots + \ket{\psi_k} \Vert)^2\\
& \leq (\Vert \ket{\psi_1}\Vert + \cdots + \Vert \ket{\psi_k} \Vert)^2\\
& \leq k \cdot(\Vert \ket{\psi_1}\Vert^2 + \cdots + \Vert \ket{\psi_k}\Vert^2).
\end{align*}
where the last step uses \Cref{eq:prop-for-real-numbers} applied to the case of $x_i = \Vert \ket{\psi_i}\Vert$.
\end{proof}

\begin{proposition}[Triangle inequality for ``$\approx_{\delta}$"]
\label{prop:triangle-inequality-for-approx_delta}
Suppose $A_1 = \{(A_1)^x_a\}, \ldots, A_{k+1} = \{(A_{k+1})^x_a\}$ is a set of matrices such that
\begin{equation*}
(A_i)^x_a \approx_{\delta_i} (A_{i+1})^x_a
\end{equation*}
for all $i \in [k]$. Then
\begin{equation*}
(A_1)^x_a \approx_{k \cdot (\delta_1 + \cdots + \delta_{k})} (A_{k+1})^x_a.
\end{equation*}
\end{proposition}
\begin{proof}
We want to bound
\begin{align*}
&~\E_\bx \sum_a \Vert ((A_1)^{\bx}_a - (A_{k+1})^{\bx}_a) \ket{\psi}\Vert^2\\
= &~\E_\bx \sum_a \Vert (((A_1)^{\bx}_a - (A_2)^{\bx}_a) + \cdots.+ ((A_{k})^{\bx}_a - (A_{k+1})^{\bx}_a))\ket{\psi}\Vert^2\\
\leq &~\E_\bx \sum_a k\cdot(\Vert ((A_1)^{\bx}_a - (A_2)^{\bx}_a) \ket{\psi}\Vert^2 + \cdots.+ \Vert((A_k)^{\bx}_a - (A_{k+1})^{\bx}_a)\ket{\psi}\Vert^2)\\
\leq &~k\cdot(\delta_1 + \cdots + \delta_{k}),
\end{align*}
where the inequality uses \Cref{prop:triangle-inequality-for-vectors-squared} applied to the vectors $((A_i)^{\bx}_a - (A_{i+1})^{\bx}_a) \ket{\psi}$ for $i \in [k]$.
\end{proof}

We note that \Cref{prop:triangle-inequality-for-approx_delta}
contrasts with the triangle inequality for ``$\approx_{\delta}$'' when applied to numbers,
i.e.\ \Cref{eq:triangle-inequality-for-numbers},
for which no multiplicative factor of~$k$ appears in the error.

The following is Fact~$4.29$ from~\cite{NW19};
however, they incorrectly claimed a final bound of $A^x_a \ot I \simeq_{\eps +\delta + \gamma} I \ot D^x_a$.
We give a new proof of this statement, albeit with a slightly weaker quantitative bound.

\begin{proposition}[Triangle inequality for ``$\simeq$"]\label{prop:simeq-triangle-inequality}
		Suppose that $A$, $B$, $C$, and~$D$ are measurements such that
		\begin{equation*}
		A^x_a \ot I \simeq_{\eps} I \ot B^x_a,
		\quad
		C^x_a \ot I \simeq_{\delta} I \ot B^x_a,
		\quad
		C^x_a \ot I \simeq_{\gamma} I \ot D^x_a.
		\end{equation*}
		Then
		\begin{equation*}
		A^x_a \ot I \simeq_{\eps + 2\sqrt{\delta + \gamma}} I \ot D^x_a.
		\end{equation*}
		\end{proposition}
\begin{proof}
Because $B$, $C$, and~$D$ are measurements, \Cref{prop:simeq-to-approx} implies that
\begin{equation*}
C^x_a \ot I \approx_{2\delta} I \ot B^x_a,
\quad
C^x_a \ot I \approx_{2\gamma} I \ot D^x_a.
\end{equation*}
The triangle inequality, \Cref{prop:triangle-inequality-for-approx_delta}, then implies that
\begin{equation*}
I \ot B^x_a \approx_{4 \delta + 4 \gamma} I \ot D^x_a.
\end{equation*}
Finally, \Cref{prop:triangle-sub} implies that
\begin{equation*}
A^x_a \ot I
\simeq_{\eps + \sqrt{4 \delta + 4 \gamma}} I \ot D^x_a.
\end{equation*}
This completes the proof.
\end{proof}

\begin{proposition}[Data processing for ``$\simeq$'']
\label{prop:simeq-data-processing}
Let $A = \{A^x_a\}$ and $B = \{B^x_a\}$ be two measurements such that
\begin{equation*}
A^x_a \ot I \simeq_{\delta} I \ot B^x_a.
\end{equation*}
Then for any function~$f$,
\begin{equation*}
A^x_{[f(a)=b]} \ot I \simeq_{\delta} I \ot B^x_{[f(a)=b]}.
\end{equation*}
\end{proposition}
\begin{proof}
We want to bound
\begin{align*}
\E_{\bx} \sum_{b \neq b'} \bra{\psi} A^{\bx}_{[f(a)=b]} \ot B^{\bx}_{[f(a)=b']}\ket{\psi}
& = \E_{\bx} \sum_{b \neq b'} \sum_{a: f(a) = b} \sum_{a': f(a') = b'} \bra{\psi} A^{\bx}_{a} \ot B^{\bx}_{a'}\ket{\psi}\\
& \leq \E_{\bx} \sum_{a \neq a'} \bra{\psi} A^{\bx}_{a} \ot B^{\bx}_{a'}\ket{\psi}\\
& \leq \delta.
\end{align*}
This completes the proof.
\end{proof}
    
The following fact is useful for translating between statements about consistency and closeness between sub-measurements.
\begin{proposition}
\label{prop:cons-sub-meas} 
Let $\{A^x_a\}$ be a sub-measurement and let $\{B^x_a\}$ be a measurement such that on average over $x$,
\[
	A^x_a \ot I \simeq_\gamma I \ot B^x_a \,.
\]
Then the following hold
	\begin{gather}
		A^x_a \ot I \approx_{\gamma} A^x_a \ot B^x_a \approx_{\gamma} A^x \ot B^x_a,
		\label{eq:closeness5}
	\end{gather}
	where $A^x = \sum_a A^x_a$.
	As a result, by \Cref{prop:triangle-inequality-for-approx_delta},
	\begin{equation*}
	A^x_a \ot I \approx_{4\gamma} A^x \ot B^x_a
	\end{equation*}
\end{proposition}
\begin{proof}
	We establish the first approximation in \Cref{eq:closeness5}:
	\begin{align*}
		&\E_{\bx} \sum_a \bra{\psi} \Big (A^{\bx}_a \ot (I - B^{\bx}_a) \Big)^2 \ket{\psi}  \\
		&\leq \E_{\bx} \sum_a \bra{\psi} A^{\bx}_a \ot (I - B^{\bx}_a) \ket{\psi} \\
		&= \E_{\bx} \sum_{\substack{a,b : \\ b \neq a}} \bra{\psi} A^{\bx}_a \ot B^{\bx}_b \ket{\psi} \\
		&\leq \gamma \,.
	\end{align*}
	The first inequality folows from the fact that $A^x_a \ot (I - B^x_a)$ has operator norm at most $1$, and the third line follows from the fact that $\{B^x_b\}$ is a complete measurement, and the last line follows from the assumption of consistency between the $A$ and $B$ (sub-)measurements.
	
	To establish the second approximation in \Cref{eq:closeness5}, we compute the difference:
	\begin{align*}
	&\E_{\bx} \sum_a \bra{\psi} \Big ( (A^{\bx} - A^{\bx}_a) \ot B^{\bx}_a \Big)^2 \ket{\psi} \\
	&\leq  \E_{\bx} \sum_a \bra{\psi} (A^{\bx} - A^{\bx}_a) \ot B^{\bx}_a  \ket{\psi} \\
	&= \E_{\bx} \sum_{\substack{a,a' : \\ a \neq a'}} \bra{\psi} A^{\bx}_{a'} \ot B^{\bx}_a  \ket{\psi} \\
	&\leq \gamma \,.
	\end{align*}
	The second line follows from the fact that  $(A^x - A^x_a) \ot B^x_a$ has operator norm at most $1$, and the last inequality follows from the consistency between the $A$ and $B$ (sub-)measurements.
\end{proof}



\begin{proposition}  \label{prop:switch-sandwich}
  Suppose $\{A^x_a\}$ is a projective sub-measurement
  satisfying
  \begin{equation}
    A^x_a \ot I \approx_{\delta} I \ot A^x_a. \label{eq:Aapproxd}
  \end{equation}
  Then for any $0 \leq B \leq I$, the following holds:
  \begin{equation}
    \E_{\bx}\sum_a \bra{\psi} A^{\bx}_a B A^{\bx}_a \ot I \ket{\psi}
    \approx_{2\sqrt{\delta}} \E_{\bx} \sum_a \bra{\psi} B
    \ot A^{\bx}_a
    \ket{\psi} \approx_{\sqrt{\delta}} \E_{\bx} \sum_a \bra{\psi}
    BA^{\bx}_a \ot I \ket{\psi} \label{eq:switch-sandwich} \end{equation}

\end{proposition}
\begin{proof}
  We will show the first approximation in \Cref{eq:switch-sandwich} in two steps. First, we show that
  \begin{equation} \label{eq:shift-right-A}
    \E_{\bx} \sum_a \bra{\psi} A^{\bx}_a B A^{\bx}_a \ot I \ket{\psi} \approx_{\sqrt{\delta}}
    \E_{\bx} \sum_a \bra{\psi} A^{\bx}_a B \ot A^{\bx}_a \ket{\psi}.
  \end{equation}
  To do so, we bound the magnitude of the difference.
  \begin{align*}
   &\Big| \E_{\bx}\sum_a \bra{\psi} (A^{\bx}_a B \otimes I)\cdot (A^{\bx}_a \ot I  - I \ot A^{\bx}_a)  \ket{\psi} \Big|\\
    \leq~&
    \Big( \E_{\bx} \sum_a \bra{\psi} A^{\bx}_a
                 B^2
                 A^{\bx}_a \otimes I \ket{\psi} \Big)^{1/2}  
                \cdot \Big(\E_{\bx} \sum_a
                 \bra{\psi} (A^{\bx}_a \ot I - I \ot A^{\bx}_a)^2
                 \ket{\psi} \Big)^{1/2} \\
               \leq~& \sqrt{\delta}.\tag{because $B \leq I$ and \eqref{eq:Aapproxd}}
  \end{align*}
  Next, we show that
  \begin{equation*}
  \eqref{eq:shift-right-A} =  \E_{\bx} \sum_a \bra{\psi} A^{\bx}_a B \ot A^{\bx}_a \ket{\psi} \approx_{\sqrt{\delta}} \E_{\bx} \sum_ a\bra{\psi} B \ot A^{\bx}_a \ket{\psi}.
  \end{equation*}
  To do so, we bound the magnitude of the difference.
  \begin{align*}
    &\Big|\E_{\bx} \sum_a \bra{\psi} (A^{\bx}_a \ot I - I \ot A^{\bx}_a) \cdot(B
    \ot A^{\bx}_a) \ket{\psi}\Big| \tag{because~$A$ is projective}\\ \leq~& \Big( \E_{\bx} \sum_a \bra{\psi}
                               (A^{\bx}_a \ot I - I \ot A^{\bx}_a)^2
                               \ket{\psi} \Big)^{1/2}  
                              \cdot \Big(\E_{\bx} \sum_a
                               \bra{\psi} B^2 \ot A^{\bx}_a \ket{\psi}
                               \Big)^{1/2} \\
                             \leq~& \sqrt{\delta}.\tag{because $B \leq I$ and \eqref{eq:Aapproxd}}
  \end{align*}
  Thus, the first approximation in \Cref{eq:switch-sandwich}
  follows. To show the second approximation, we bound the magnitude of the difference.
  \begin{align*}
  &\Big|\E_{\bx}\sum_a \bra{\psi} (B\otimes I)  \cdot ( A^{\bx}_a \ot I - I \ot A^{\bx}_a )
    \ket{\psi}\Big|\\
    =~&\Big|\E_{\bx} \bra{\psi} (B\otimes I)\cdot \sum_a ( A^{\bx}_a \ot I - I \ot A^{\bx}_a )
    \ket{\psi}\Big|\\
    \leq~& \Big( \E_{\bx} \bra{\psi} B^2 \otimes I \ket{\psi} \Big)^{1/2}
    \cdot \Big( \E_{\bx} \sum_{a,b} \bra{\psi}
    (A^{\bx}_a \ot I - I \ot A^{\bx}_a) \cdot(A^{\bx}_b \ot I
    - I \ot A^{\bx}_b) \ket{\psi} \Big)^{1/2}.
  \end{align*}
  The first term in the product is at most $1$ because $B \leq I$.
  We bound the expression inside the second square root as follows.
  \begin{align*}
    &\E_{\bx} \sum_{a,b} \bra{\psi}    (A^{\bx}_a \ot I - I \ot A^{\bx}_a) \cdot(A^{\bx}_b \ot I
    - I \ot A^{\bx}_b) \ket{\psi}  \\
    =~&\E_{\bx} \sum_{a, b} \bra{\psi}(A^{\bx}_a A^{\bx}_b \ot I + I \otimes A^{\bx}_a A^{\bx}_b - A^{\bx}_a \otimes A^{\bx}_b - A^{\bx}_b \otimes A^{\bx}_a)\ket{\psi}\\
    =~& \E_{\bx}\sum_a
                                                \bra{\psi} (A^{\bx}_a)^2
                                                \ot I \ket{\psi} +
                                                 \E_{\bx}\sum_a
                                                \bra{\psi} I \ot (A^{\bx}_a)^2 \ket{\psi} -
                                                 2\cdot\E_{\bx} \sum_{a,b} \bra{\psi}A^{\bx}_a
                                                \ot A^{\bx}_b
                                                \ket{\psi} \tag{because~$A$ is projective}
                                                \\
    \leq~& \E_{\bx}\sum_a
                                                \bra{\psi} (A^{\bx}_a)^2
                                                \ot I \ket{\psi} +
                                                 \E_{\bx}\sum_a
                                                \bra{\psi} I \ot (A^{\bx}_a)^2 \ket{\psi} -
                                                 2\cdot\E_{\bx} \sum_{a} \bra{\psi}A^{\bx}_a
                                                \ot A^{\bx}_a
                                                \ket{\psi}
                                                \\
       =~& \E_{\bx}\sum_a
                                                \bra{\psi} (A^{\bx}_a
                                                \ot I  - I \ot A^{\bx}_a)^2
                                                \ket{\psi}
       \\
   \leq~& \delta,
  \end{align*}
  where the last step uses \Cref{eq:Aapproxd}.
\end{proof}

\begin{proposition}\label{prop:completeness-transfer-projective-P}
Let $A = \{A^x_a\}$ be a sub-measurement 
and let $P$ be a projective sub-measurement such that $A^x_a \ot I \approx_{\eps} P_a^x \ot I$.
Then
\begin{equation*}
\bra{\psi} A \ot I \ket{\psi}
\geq  \bra{\psi} P \ot I \ket{\psi} - 2\sqrt{\eps}.
\end{equation*}
\end{proposition}
\begin{proof}
We calculate:
\begin{align*}
\bra{\psi} P \ot I \ket{\psi}
&= \E_{\bx} \sum_a \bra{\psi} P^{\bx}_a \ot I \ket{\psi} \\
&= \E_{\bx} \sum_a \bra{\psi} (P^{\bx}_a)^2 \ot I \ket{\psi} \tag{because~$P$ is projective}\\
&\approx_{\sqrt{\eps}} \E_{\bx} \sum_a \bra{\psi} (A^{\bx}_a \cdot P^{\bx}_a) \ot I \ket{\psi}
									\tag{by \Cref{prop:easy-approx-from-approx-delta}} \\
&\approx_{\sqrt{\eps}} \E_{\bx} \sum_a \bra{\psi} (A^{\bx}_a)^2 \ot I \ket{\psi}
									\tag{by \Cref{prop:easy-approx-from-approx-delta}} \\
&\leq \E_{\bx} \sum_a \bra{\psi} A^{\bx}_a \ot I \ket{\psi}\\
& = \bra{\psi} A \ot I \ket{\psi}.
\end{align*}
This completes the proof.
\end{proof}

\subsubsection{Strong self-consistency}

An important property of a sub-measurement~$A = \{A^x_a\}$
is that if both provers measure using~$A$,
then they receive the same outcome.
It seems natural to study this using \emph{self-consistency} of~$A$, i.e. the number~$\delta$ such that
\begin{equation*}
A^x_a \ot I \simeq_{\delta} I \ot A^x_a.
\end{equation*}
However, when~$A$ is a sub-measurement,
being $\delta$-self consistent
only implies the following weaker condition:
if both provers measure using~$A$ and one of them receives~$\ba$,
then the other will most likely either receive~$\ba$ \emph{or not receive any outcome whatsoever}.
This motivates defining the following stronger notion of self-consistency.

\begin{definition}[Strong self consistency]\label{def:strong-self-consistency}
Let $\ket{\psi}$ be a permutation-invariant state in $\calH \ot \calH$
and let $A = \{A^x_a\}$ be a sub-measurement acting on~$\calH$.
Then $A$ is \emph{$\delta$-strongly self consistent} if
\begin{equation*}
\E_{\bx} \sum_a \bra{\psi} A^{\bx}_a \otimes A^{\bx}_a \ket{\psi}
\geq
\bra{\psi} A \otimes I \ket{\psi} -  \delta.
\end{equation*}
\end{definition}

We now relate strong self-consistency to our two notions of similarity.
First, we show that strong self-consistency is indeed
a stronger condition than $A^x_a \ot I \simeq_{\delta} I \ot A^x_a$,
at least for sub-measurements.

\begin{proposition}\label{prop:other-two-notions-of-self-consistency}
Let $\ket{\psi}$ be a permutation-invariant state,
and let $A = \{A^x_a\}$ be a sub-measurement. If
\begin{equation*}
\E_{\bx} \sum_a \bra{\psi} A^{\bx}_a \otimes A^{\bx}_a \ket{\psi} \geq \bra{\psi} A \otimes I \ket{\psi} - \delta
\end{equation*}
then $A^x_a \otimes I \simeq_{\delta} I \otimes A^x_a$.
This is an ``if and only if" if $A$ is a measurement.
\end{proposition}
\begin{proof}
For a sub-measurement~$A$,
\begin{align*}
&\E_{\bx} \sum_{a \neq b} \bra{\psi} A^{\bx}_a \ot A^{\bx}_b \ket{\psi}\\
=~& \E_{\bx} \sum_{a} \bra{\psi} A^{\bx}_a \ot (A^{\bx}- A^{\bx}_a) \ket{\psi}\\
\leq~&\E_{\bx} \sum_{a} \bra{\psi} A^{\bx}_a \ot (I- A^{\bx}_a) \ket{\psi}\\
=~&\E_{\bx} \sum_{a} \bra{\psi} A^{\bx}_a \ot I \ket{\psi}
		- \E_{\bx} \sum_{a} \bra{\psi} A^{\bx}_a \ot A^{\bx}_a \ket{\psi}\\
=~& \bra{\psi} A \ot I \ket{\psi}
		- \E_{\bx} \sum_{a} \bra{\psi} A^{\bx}_a \ot A^{\bx}_a \ket{\psi}.
\end{align*}
This is at most~$\delta$ if $A$ is $\delta$-strongly self-consistent.
On the other hand, if~$A$ is a measurement, then the inequality becomes an equality.
Hence, if $A^x_a \ot I \simeq_{\delta} I \otimes A^x_a$,
then $A$ is $\delta$-strongly self-consistent.
\end{proof}

Next, we show that strong self-consistency is also a stronger 
condition than $A^x_a \ot I \approx_{2\delta} I \ot A^x_a$, at least for non-projective measurements.

\begin{proposition}\label{prop:two-notions-of-self-consistency}
Let $\ket{\psi}$ be a permutation-invariant state,
and let $A = \{A^x_a\}$ be a sub-measurement. If
\begin{equation*}
\E_{\bx} \sum_a \bra{\psi} A^{\bx}_a \otimes A^{\bx}_a \ket{\psi} \geq \bra{\psi} A \otimes I \ket{\psi} - \delta
\end{equation*}
then $A^x_a \otimes I \approx_{2\delta} I \otimes A^x_a$.
This is an ``if and only if" if $A$ is projective.
\end{proposition}
\begin{proof}
For general (i.e.\ not necessarily projective)~$A$
\begin{align}
&\E_{\bx} \sum_a \Vert(A^{\bx}_a \otimes I - I \otimes A^{\bx}_a) \ket{\psi}\Vert^2\nonumber\\
=~&\E_{\bx} \sum_a \bra{\psi} (A^{\bx}_a \otimes I - I \otimes A^{\bx}_a)^2 \ket{\psi}\nonumber\\
=~&2\cdot \Big( \E_{\bx} \sum_a \bra{\psi} (A^{\bx}_a)^2 \otimes I \ket{\psi} -  \E_{\bx} \sum_a \bra{\psi} A^{\bx}_a \otimes A^{\bx}_a \ket{\psi}\Big)\nonumber\\
\leq~&2\cdot \Big( \E_{\bx} \sum_a \bra{\psi} A^{\bx}_a \otimes I \ket{\psi} -  \E_{\bx} \sum_a \bra{\psi} A^{\bx}_a \otimes A^{\bx}_a \ket{\psi}\Big).\label{eq:here's-where-projectivity-would-help}
\end{align}
This is at most~$2\cdot \delta$ if
\begin{equation*}
\E_{\bx} \sum_a \bra{\psi} A^{\bx}_a \otimes A^{\bx}_a \ket{\psi} \geq \E_{\bx} \sum_a \bra{\psi} A^{\bx}_a \otimes I \ket{\psi} - \delta.
\end{equation*}
If~$A$ is projective, then \Cref{eq:here's-where-projectivity-would-help} becomes an equality,
and so $A^x_a \otimes I \approx_{2\delta} I \otimes A^x_a$ implies that
\begin{equation*}
 \E_{\bx} \sum_a \bra{\psi} A^{\bx}_a \otimes I \ket{\psi} -  \E_{\bx} \sum_a \bra{\psi} A^{\bx}_a \otimes A^{\bx}_a \ket{\psi}
 = \frac{1}{2} \cdot \left(\E_{\bx} \sum_a \Vert(A^{\bx}_a \otimes I - I \otimes A^{\bx}_a) \ket{\psi}\Vert^2\right) \leq \delta.\qedhere
\end{equation*}
\end{proof}

Hence, we may also refer to the condition $A^x_a \ot I \approx_{2\delta} I \ot A^x_a$
as ``strong self-consistency'' if~$A$ is projective.

For the remainder of the section, we will prove various properties of strongly self-consistent sub-measurements.

\begin{proposition}\label{prop:two-notions-of-self-consistency-after-evaluation}
Let $\ket{\psi}$ be a permutation-invariant state,
and let $A = \{A^x_a\}$ be a sub-measurement such that
\begin{equation*}
\E_{\bx} \sum_a \bra{\psi} A^{\bx}_a \otimes A^{\bx}_a \ket{\psi} \geq  \bra{\psi} A \otimes I \ket{\psi} - \delta.
\end{equation*}
Then for any function~$f$,
$A^x_{[f(a)=b]} \otimes I \approx_{2\delta} I \otimes A^x_{[f(a)=b]}$.
\end{proposition}
\begin{proof}
Our goal is to bound
\begin{align}
&\E_{\bx} \sum_b \Vert(A^{\bx}_{[f(a)=b]} \otimes I - I \otimes A^{\bx}_{[f(a)=b]}) \ket{\psi}\Vert^2\nonumber\\
=~&\E_{\bx} \sum_b \bra{\psi} (A^{\bx}_{[f(a)=b]} \otimes I - I \otimes A^{\bx}_{[f(a)=b]})^2 \ket{\psi}\nonumber\\
=~&2\cdot \Big( \E_{\bx} \sum_b \bra{\psi} (A^{\bx}_{[f(a)=b]})^2 \otimes I \ket{\psi} -  \E_{\bx} \sum_b \bra{\psi} A^{\bx}_{[f(a)=b]} \otimes A^{\bx}_{[f(a)=b]} \ket{\psi}\Big)\nonumber\\
\leq~&2\cdot \Big( \E_{\bx} \sum_b \bra{\psi} A^{\bx}_{[f(a)=b]} \otimes I \ket{\psi} -  \E_{\bx} \sum_b \bra{\psi} A^{\bx}_{[f(a)=b]} \otimes A^{\bx}_{[f(a)=b]} \ket{\psi}\Big).\nonumber\\
=~&2\cdot \Big(  \bra{\psi} A \otimes I \ket{\psi} -  \E_{\bx} \sum_b \bra{\psi} A^{\bx}_{[f(a)=b]} \otimes A^{\bx}_{[f(a)=b]} \ket{\psi}\Big).\label{eq:finishing-this-up}
\end{align}
The second term in \Cref{eq:finishing-this-up} can be bounded by
\begin{align*}
 \E_{\bx} \sum_b \bra{\psi} A^{\bx}_{[f(a)=b]} \otimes A^{\bx}_{[f(a)=b]} \ket{\psi}
 &=  \E_{\bx} \sum_{a, a': f(a) = f(a')} \bra{\psi} A^{\bx}_{a} \otimes A^{\bx}_{a'} \ket{\psi}\\
 &\geq\E_{\bx} \sum_{a} \bra{\psi} A^{\bx}_{a} \otimes A^{\bx}_{a} \ket{\psi}\\
 &\geq \bra{\psi} A \ot I \ket{\psi} -\delta.
\end{align*}
Hence,
\begin{equation*}
\eqref{eq:finishing-this-up} \leq 
	2\cdot \Big( \bra{\psi} A \otimes I \ket{\psi} - (\bra{\psi} A \ot I \ket{\psi} -\delta)\Big)
= 2\delta.
\end{equation*}
This concludes the proof.
\end{proof}




\begin{proposition}\label{prop:completeness-transfer-self-consistent-A}
Let $\ket{\psi}$ be a permutation-invariant state,
and let $A = \{A^x_a\}$ be a sub-measurement such that
\begin{equation*}
\E_{\bx} \sum_a \bra{\psi} A^{\bx}_a \otimes A^{\bx}_a \ket{\psi} \geq  \bra{\psi} A \otimes I \ket{\psi} - \delta.
\end{equation*}
In addition, let $B$ be a sub-measurement such that $A^x_a \ot I \approx_{\eps} B_a^x \ot I$.
Then
\begin{equation*}
\bra{\psi} B \ot I \ket{\psi}
\geq \bra{\psi} A \ot I \ket{\psi} - \delta - 2\sqrt{\eps}.
\end{equation*}
\end{proposition}
\begin{proof}
We calculate:
\begin{align*}
\bra{\psi} B \ot I \ket{\psi}
&= \E_{\bx} \sum_a \bra{\psi} B^{\bx}_a \ot I \ket{\psi} \\
&\geq \E_{\bx} \sum_a \bra{\psi} B^{\bx}_a \ot B^{\bx}_a \ket{\psi} \\
& \geq \E_{\bx} \sum_a \bra{\psi} A^{\bx}_a \ot B^{\bx}_a \ket{\psi} - \sqrt{\eps}
				\tag{by \Cref{prop:easy-approx-from-approx-delta}} \\
& \geq \E_{\bx} \sum_a \bra{\psi} A^{\bx}_a \ot A^{\bx}_a \ket{\psi} - 2\sqrt{\eps}
				\tag{by \Cref{prop:easy-approx-from-approx-delta}} \\
& \geq  \bra{\psi} A \ot I \ket{\psi} - \delta - 2\sqrt{\eps}.
\end{align*}
This completes the proof.
\end{proof}


\begin{proposition}\label{prop:self-consistency-implies-data-processing}
Let $\ket{\psi}$ be a permutation-invariant state,
and let $A = \{A^x_a\}$ be a sub-measurement such that
\begin{equation*}
\E_{\bx} \sum_a \bra{\psi} A^{\bx}_a \otimes A^{\bx}_a \ket{\psi} \geq  \bra{\psi} A \otimes I \ket{\psi} - \delta.
\end{equation*}
In addition, let $P$ be a projective sub-measurement such that $P^x_a \ot I \approx_{\eps} A_a^x \ot I$.
Then for any function~$f$,
\begin{equation*}
P^x_{[f(a)=b]} \ot I \approx_{8\delta + 8\sqrt{\eps}} A_{[f(a)=b]}^x \ot I.
\end{equation*}
\end{proposition}
\begin{proof}
We will begin by showing that
\begin{equation}\label{eq;what-we-want-to-prove-but-on-wrong-side}
P^x_{[f(a)=b]} \ot I \approx_{2\delta + 4\sqrt{\eps}} I \ot A_{[f(a)=b]}^x.
\end{equation}
To do so, our goal is to bound
\begin{align}
&\E_{\bx} \sum_b \Vert(P^{\bx}_{[f(a)=b]} \ot I - I \ot A_{[f(a)=b]}^{\bx}) \ket{\psi} \Vert^2\nonumber\\
=~&\E_{\bx} \sum_b \bra{\psi} (P^{\bx}_{[f(a)=b]}  \ot I -  I \ot A_{[f(a)=b]}^{\bx})^2 \ket{\psi} \nonumber\\
=~&\E_{\bx} \sum_b \bra{\psi} \Big((P^{\bx}_{[f(a)=b]})^2 \ot I 
		+ I \ot (A^{\bx}_{[f(a)=b]})^2 
		- 2\cdot P^{\bx}_{[f(a)=b]} \ot A^{\bx}_{[f(a)=b]}\Big) \ket{\psi}\nonumber\\
\leq~&\E_{\bx} \sum_b \bra{\psi} \Big(P^{\bx}_{[f(a)=b]} \ot I 
		+ I \ot A^{\bx}_{[f(a)=b]}
		- 2\cdot P^{\bx}_{[f(a)=b]} \ot A^{\bx}_{[f(a)=b]}\Big) \ket{\psi}\nonumber\\
=~&\bra{\psi} P \ot I \ket{\psi} + \bra{\psi} I \ot A \ket{\psi}
		- 2\cdot \E_{\bx} \sum_b \bra{\psi}  P^{\bx}_{[f(a)=b]} \ot A^{\bx}_{[f(a)=b]} \ket{\psi}.
				\label{eq:gonna-handle-third-term}
\end{align}
By \Cref{prop:completeness-transfer-projective-P}, the first term in \Cref{eq:gonna-handle-third-term} is at most
\begin{equation*}
\bra{\psi} P \ot I \ket{\psi}
\leq \bra{\psi} A \ot I \ket{\psi} + 2\sqrt{\eps}.
\end{equation*}
As for the third term in \Cref{eq:gonna-handle-third-term}, we can bound it by
\begin{align*}
\E_{\bx} \sum_b \bra{\psi}  P^{\bx}_{[f(a)=b]} \ot A^{\bx}_{[f(a)=b]} \ket{\psi}
&= \E_{\bx} \sum_{a, a': f(a) = f(a')} \bra{\psi}  P^{\bx}_{a} \ot A^{\bx}_{a'} \ket{\psi}\\
&\geq \E_{\bx} \sum_{a} \bra{\psi}  P^{\bx}_{a} \ot A^{\bx}_{a} \ket{\psi}\\
&\geq \E_{\bx} \sum_{a} \bra{\psi}  A^{\bx}_{a} \ot A^{\bx}_{a} \ket{\psi} - \sqrt{\eps}
						\tag{by \Cref{prop:easy-approx-from-approx-delta}} \\
&\geq \bra{\psi} A \otimes I \ket{\psi} - \delta - \sqrt{\eps}.
\end{align*}
Putting everything together,
\begin{equation*}
\eqref{eq:gonna-handle-third-term}
\leq (\bra{\psi} A \ot I \ket{\psi} + 2\sqrt{\eps}) + \bra{\psi} A \ot I \ket{\psi}
	- 2 \cdot(\bra{\psi} A \otimes I \ket{\psi} - \delta - \sqrt{\eps})
= 2\delta + 4\sqrt{\eps}.
\end{equation*}
This proves \Cref{eq;what-we-want-to-prove-but-on-wrong-side}.


\Cref{prop:two-notions-of-self-consistency-after-evaluation} implies that
\begin{equation*}
A^x_{[f(a)=b]} \ot I \approx_{2\delta} I \ot A^x_{[f(a)=b]}.
\end{equation*}
Hence, by \Cref{eq;what-we-want-to-prove-but-on-wrong-side},
\begin{equation*}
P^x_{[f(a)=b]} \ot I
\approx_{2\delta + 4\sqrt{\eps}} I \ot A_{[f(a)=b]}^x
\approx_{2\delta} A^x_{[f(a)=b]} \ot I.
\end{equation*}
Thus, by \Cref{prop:triangle-inequality-for-approx_delta},
\begin{equation*}
P^x_{[f(a)=b]} \ot I
\approx_{8\delta + 8\sqrt{\eps}}A^x_{[f(a)=b]} \ot I.
\end{equation*}
This concludes the proof.
\end{proof}

\begin{proposition}\label{prop:completing-to-measurement}
		Let $\ket{\psi}$ be a permutation-invariant state,
		and let $A = \{A_a\}$ be a measurement such that
		\begin{equation*}
		\E_{\bx} \sum_a \bra{\psi} A_a \otimes A_a \ket{\psi} \geq  \bra{\psi} A \otimes I \ket{\psi} - \zeta.
		\end{equation*}
		Suppose $B = \{B_a\}$ is a sub-measurement such that
		$A_a\ot I \approx_{\delta} B_a\ot I$.
		Let $C = \{C_a\}$ be a measurement in which there is an~$a^*$
		such that $C_{a^*} = B_{a^*} + (I-B)$ and $C_a = B_a$ for all $a \neq a^*$.
		Then $A_a \ot I \approx_{2\delta + 4 \sqrt{\delta} + 2\zeta} C_a \ot I$.
		\end{proposition}
		
		Before proving this, we need the following proposition.
		\begin{proposition}\label{prop:cool-prop}
		Let $\ket{\psi}$ be a permutation-invariant state,
		and let $A = \{A_a\}$ be a sub-measurement such that
		\begin{equation*}
		\E_{\bx} \sum_a \bra{\psi} A_a \otimes A_a \ket{\psi} \geq  \bra{\psi} A \otimes I \ket{\psi} - \zeta.
		\end{equation*}
		Then
		\begin{equation*}
		\sum_a \bra{\psi} (A_a)^2 \otimes I \ket{\psi} \geq \sum_a \bra{\psi} A_a \otimes I \ket{\psi} -\zeta.
		\end{equation*}
		\end{proposition}
		\begin{proof}
		Applying Cauchy-Schwarz, we have
\begin{align}
\sum_a \bra{\psi} A_a \otimes A_a \ket{\psi}
&= \sum_a \bra{\psi} (A_a \ot I) \cdot (I \ot A_a) \ket{\psi}\nonumber\\
&\leq \sqrt{\sum_a \bra{\psi} (A_a)^2 \ot I \ket{\psi}} \cdot \sqrt{\sum_a \bra{\psi} I \ot (A_a)^2 \ket{\psi}}\nonumber\\
& = \sum_a \bra{\psi} (A_a)^2 \ot I \ket{\psi}.\label{eq:As-on-same-side-rewrite}
\end{align}
As a result,
\begin{align*}
\sum_a \bra{\psi} (A_a)^2 \otimes I \ket{\psi}
&\geq \sum_a \bra{\psi} A_a \otimes A_a \ket{\psi}\tag{by \Cref{eq:As-on-same-side-rewrite}}\\
&\geq \sum_a \bra{\psi} A_a \otimes I \ket{\psi} -\zeta. \tag{by self-consistency of~$A$}
\end{align*}
This completes the proof.
\end{proof}
		Now we prove \Cref{prop:completing-to-measurement}.
		\begin{proof}[Proof of \Cref{prop:completing-to-measurement}]
		By \Cref{prop:triangle-inequality-for-vectors-squared} (the triangle inequality for vectors squared),
		\begin{align*}
		&\phantom{=}\sum_a \Vert(A_a - C_a)\ot I \ket{\psi} \Vert^2\\
		& = \sum_{a \neq a^*} \Vert(A_a - B_a)\ot I \ket{\psi} \Vert^2 + \Vert(A_{a^*}  - (I-B + B_{a^*}))\ot I\ket{\psi}\Vert^2\\
		& \leq \sum_{a \neq a^*} \Vert(A_a - B_a)\ot I \ket{\psi} \Vert^2 + 2 \cdot \Vert(A_{a^*}  -  B_{a^*})\ot I\ket{\psi}\Vert^2 + 2 \cdot \Vert(I - B) \ot I\ket{\psi}\Vert^2\\
		& \leq 2\sum_{a \neq a^*} \Vert(A_a - B_a)\ot I \ket{\psi} \Vert^2 + 2 \cdot \Vert(A_{a^*}  -  B_{a^*})\ot I\ket{\psi}\Vert^2 + 2 \cdot \Vert(I - B) \ot I\ket{\psi}\Vert^2\\
		& = 2\sum_{a} \Vert(A_a - B_a)\ot I \ket{\psi} \Vert^2 + 2 \cdot \Vert(I - B)\ot I \ket{\psi}\Vert^2\\
		& \leq 2 \delta + 2 \cdot \Vert(I - B)\ot I \ket{\psi}\Vert^2. \tag{because $A_a \ot I \approx_{\delta} B_a \ot I$}
		\end{align*}
		The second term we can bound as follows.
		\begin{align}
		\Vert(I - B) \ot I \ket{\psi}\Vert^2
		& = \bra{\psi} (I-B)^2 \ot I \ket{\psi}\nonumber\\
		& \leq \bra{\psi} (I-B)\ot I \ket{\psi} \tag{because~$B$ is a sub-measurement}\\
		& = 1 - \bra{\psi} B\ot I \ket{\psi}\nonumber\\
		& = 1 - \sum_a \bra{\psi} B_a \ot I \ket{\psi}\nonumber\\
		& \leq 1 -\sum_a \bra{\psi} B_a^2 \ot I \ket{\psi}.\label{eq:to-return-later-whatevs}
		\end{align}
		Now, by \Cref{prop:easy-approx-from-approx-delta} and the fact that $A_a\ot I \approx_{\delta} B_a\ot I$,
		\begin{align*}
		\sum_a \bra{\psi} B_a^2 \ot I \ket{\psi}
		& \approx_{\sqrt{\delta}} \sum_a \bra{\psi} (A_a \cdot B_a)\ot I \ket{\psi}\\
		& \approx_{\sqrt{\delta}} \sum_a \bra{\psi} A_a^2\ot I \ket{\psi}\\
		& \geq  \bra{\psi} A \ot I \ket{\psi}- \zeta \tag{by \Cref{prop:cool-prop}}\\
		& = 1 - \zeta. \tag{because~$A$ is a measurement}
		\end{align*}
		Hence, by \Cref{eq:to-return-later-whatevs},
		\begin{equation*}
		\Vert(I - B) \ot I \ket{\psi}\Vert^2 \leq 1 -\sum_a \bra{\psi} B_a^2 \ot I \ket{\psi} \leq 2\sqrt{\delta} + \zeta.
		\end{equation*}
		This concludes the proof.
		\end{proof}


